%% =====================================================================
%% Journal of Science and Technology (KNUST) -- revised submission (R2)
%% Opoku, Adu, Emmanuel & Owusu
%% Compile with pdflatex (twice). No external .bib or figure files needed.
%% All tables are generated by make_tables.R from the R replication output (Supplementary S1).
%% \showverifytrue prints red author-action flags; set \showverifyfalse
%% for the final upload once they are resolved.
%% =====================================================================
\documentclass[11pt,a4paper]{article}

\usepackage[T1]{fontenc}
\usepackage[utf8]{inputenc}
\usepackage{mathptmx}
\usepackage[scaled=0.92]{helvet}
\usepackage{courier}
\usepackage[margin=2.5cm]{geometry}
\usepackage{amsmath,amssymb,bm}
\usepackage{booktabs,array,threeparttable}
\usepackage[table]{xcolor}
\usepackage{pgfplots}
\pgfplotsset{compat=1.17}
\usepgfplotslibrary{groupplots}
\usepackage{setspace}
\usepackage[authoryear,round]{natbib}
\usepackage{authblk}
\usepackage{lineno}
\usepackage{caption}
\captionsetup{font=small,labelfont=bf,skip=4pt}
\usepackage{fancyhdr}
\usepackage[hidelinks]{hyperref}
\usepackage{doi}

\newif\ifshowverify
\showverifytrue
\newcommand{\verify}[1]{\ifshowverify\textcolor{red}{\,[\textit{Authors to complete:} #1]}\fi}

\newcommand{\COtwo}{CO\textsubscript{2}}
\newcommand{\E}{\mathbb{E}}

\pagestyle{fancy}
\fancyhf{}
\fancyhead[L]{\small\itshape Estimator performance in unbalanced climate panels}
\fancyhead[R]{\small\itshape Opoku et al.}
\fancyfoot[C]{\small\thepage}
\renewcommand{\headrulewidth}{0.4pt}

\onehalfspacing
\linenumbers

\title{\vspace{-1.5em}\textbf{Evaluating Estimator Performance in Unbalanced Dynamic Climate--Economy Panels: A Monte Carlo Comparison of Fixed Effects, Instrumental Variables, Difference-GMM and System-GMM}}

\author[1]{Emmanuel Opoku}
\author[1]{Alfred Asiwome Adu}
\author[1]{Harris Emmanuel}
\author[1]{Owusu Enock\thanks{Corresponding author. E-mail: \texttt{enockowusu52@gmail.com}}}
\affil[1]{Department of Statistics and Actuarial Science, Kwame Nkrumah University of Science and Technology (KNUST), Kumasi, Ghana}
\date{}

\begin{document}
\maketitle
\thispagestyle{fancy}

%% =====================================================================
\begin{abstract}
\noindent
Climate--economy studies often estimate dynamic panel models in which output is persistent, emissions are jointly determined with output, and country coverage is incomplete. This study compares the finite-sample performance of the fixed-effects (FE), Anderson--Hsiao instrumental-variables (IV), Difference-GMM and System-GMM estimators under these conditions. A fully specified Monte Carlo design generates 1{,}000 replications of panels that begin with 135 countries and 50 years, made unbalanced by staggered entry and random deletion. GDP persistence and the \COtwo{} elasticity are both 0.70, and the share of \COtwo{} innovations transmitted to the output disturbance is set at $\lambda\in\{0,0.3,0.6,0.9\}$, where $\lambda=0$ is an exogenous baseline. Estimators are assessed by bias, standard deviation, root mean squared error (RMSE), median absolute error and confidence-interval coverage, alongside the frequency of failed fits, instrument counts and replication-level specification tests. With exogenous emissions, FE was the most accurate estimator. Once emissions were endogenous, FE bias grew with $\lambda$, reaching $-34\%$ for persistence and $+106\%$ for the \COtwo{} elasticity at $\lambda=0.9$. The exactly identified IV estimator suffered from weak instruments and extreme outliers, and Difference-GMM estimated the \COtwo{} elasticity imprecisely and with increasing upward bias. System-GMM, whose additional moment conditions hold by construction in this design, had the lowest RMSE for both coefficients whenever $\lambda>0$, reducing the \COtwo{} RMSE at $\lambda=0.9$ by 86.4\% relative to Difference-GMM and 83.3\% relative to FE. Its confidence intervals nonetheless under-covered (88--93\%), and its two-step Hansen test over-rejected valid instruments. The findings support System-GMM only under conditions resembling those simulated, chiefly endogenous persistent regressors and mean-stationary initial conditions.
\end{abstract}

\noindent\textbf{Keywords:} Unbalanced dynamic panels; Monte Carlo simulation; endogeneity; Difference-GMM; System-GMM; climate econometrics.

\noindent\textbf{JEL classification:} C15; C23; C26; Q54.

%% =====================================================================
\section{Introduction}

Climate change has become a central concern in empirical economics because of its implications for growth, agricultural productivity and long-term development. \citet{Dell2012} showed that higher temperatures reduce economic growth in poorer countries, and \citet{Burke2015} documented a non-linear global relationship between temperature and economic production. \citet{Auffhammer2013} and \citet{Hsiang2016} stressed that credible evidence in this field depends on careful measurement and econometric identification. The choice of estimator can therefore determine the sign and size of the climate--economy relationships that inform policy.

Panel-data methods dominate this literature because they exploit variation across countries and over time while absorbing unobserved country characteristics. Climate--economy panels nevertheless combine features that are individually well understood but rarely studied together. Output is highly persistent; climatic series exhibit short- and long-range serial dependence \citep{Gao2024} and common global components; observation records are frequently incomplete \citep{Atak2011}; and carbon dioxide (\COtwo) emissions are plausibly determined jointly with gross domestic product (GDP), so that they are endogenous when they enter the GDP equation.

Each standard estimator handles only part of this problem. The fixed-effects (FE) estimator removes time-invariant country effects, but in a dynamic model the within-transformed lagged dependent variable remains correlated with the transformed disturbance \citep{Nickell1981}, and the bias can remain material for moderately long panels \citep{Judson1999}. The instrumental-variables (IV) approach of \citet{Anderson1981,Anderson1982} first-differences the model and instruments the differenced lagged dependent variable with lagged levels, but such instruments become weak as persistence rises. Difference-GMM \citep{Arellano1991} uses lagged levels as instruments for the differenced equation, and System-GMM \citep{Arellano1995,Blundell1998} adds level equations instrumented by lagged differences, which strengthens identification under persistence at the cost of additional initial-condition restrictions \citep{Ullah2018}. System-GMM is not immune to weak instruments \citep{Bun2010} or finite-sample bias \citep{Hayakawa2007}, and unrestricted instrument counts grow quadratically in $T$, overfitting endogenous regressors and weakening specification tests \citep{Roodman2009a,Roodman2009b}.

Panel imbalance adds a further difficulty, because staggered entry and intermittent gaps reduce the availability of lags, differences and instruments. \citet{Albarran2019} developed GMM methods for incomplete and rotating panels, \citet{Kiviet1995} and \citet{Judson1999} examined the finite-sample behaviour of dynamic-panel estimators, and \citet{Bruno2005} extended bias-corrected FE estimation to unbalanced panels. None of these studies answers the applied question posed here: \emph{how do FE, IV, Difference-GMM and System-GMM compare within a single climate--economy data-generating process in which persistence, endogenous emissions, common climatic shocks and a specified missing-observation mechanism coexist?}

This study addresses that question with a fully specified and reproducible Monte Carlo design. The contribution is threefold. First, the four estimators are compared under four endogeneity settings, including an exogenous baseline ($\lambda=0$) that separates the bias caused by dynamic structure from the bias caused by endogenous emissions. Second, the study validates the simulation itself: it reports retained samples, failed fits, instrument counts, the correlation of every instrument with the true disturbance, and the behaviour of the usual specification tests across all 1{,}000 replications rather than in a single draw. Third, it translates the results into conditional guidance for applied work. The contribution is a finite-sample comparison under a stated design, not a new estimator, and no estimator is presumed to dominate for every coefficient or setting.

Section~\ref{sec:methods} sets out the data-generating process, estimators and performance measures. Section~\ref{sec:results} reports the validation evidence, an illustrative replication and the Monte Carlo results. Section~\ref{sec:discussion} discusses them, and Section~\ref{sec:conclusion} concludes.

%% =====================================================================
\section{Materials and Methods}\label{sec:methods}

\subsection{Structural model}

For country $i=1,\dots,N$ and year $t=1,\dots,T$, log GDP follows
\begin{equation}\label{eq:model}
y_{it}=\alpha y_{i,t-1}+\beta_1 q_{it}+\beta_2 p_{it}+\beta_3 h_{it}+\mu_i+\varepsilon_{it},
\end{equation}
where $q_{it}$ is log \COtwo{} emissions, $p_{it}$ is log precipitation, $h_{it}$ is the temperature anomaly (\textdegree C), $\mu_i$ is a time-invariant country effect and $\varepsilon_{it}$ is an idiosyncratic disturbance. The true parameters are
\begin{equation}\label{eq:theta}
\bm\theta_0=(\alpha,\beta_1,\beta_2,\beta_3)'=(0.70,\;0.70,\;0.003,\;-0.015)'.
\end{equation}
$\beta_1$ and $\beta_2$ are short-run elasticities, and $\beta_3$ is a semi-elasticity: a one-degree rise in the anomaly lowers GDP by about 1.5\%, other regressors held fixed. The implied long-run \COtwo{} elasticity is $\beta_1/(1-\alpha)\approx2.33$. The calibration is illustrative and is not intended to reproduce estimates from any particular empirical study.

\subsection{Data-generating process}\label{sec:dgp}

All stochastic components are mutually independent Gaussian draws unless stated otherwise. Country heterogeneity is generated from three independent standard normal variables $z_i$, $w_i$ and $v_i$:
\begin{align}
\mu_i &= \sigma_\mu z_i, \label{eq:mu}\\
m^q_i &= \bar m_q+\sigma_{mq}\big(\varrho z_i+\sqrt{1-\varrho^2}\,w_i\big), \label{eq:mq}\\
m^p_i &= \bar m_p+\sigma_{mp}v_i. \label{eq:mp}
\end{align}
The country mean of log emissions, $m^q_i$, is correlated with the country effect ($\varrho=0.6$), so that the regressors are correlated with $\mu_i$ and random-effects assumptions fail, as they plausibly do in practice. The regressors evolve as
\begin{align}
q_{it} &= m^q_i+\rho_q\,(q_{i,t-1}-m^q_i)+\eta_{it}, & \eta_{it}&\sim N(0,\sigma_\eta^2), \label{eq:q}\\
p_{it} &= m^p_i+\rho_p\,(p_{i,t-1}-m^p_i)+\zeta_{it}, & \zeta_{it}&\sim N(0,\sigma_\zeta^2), \label{eq:p}\\
h_{it} &= \rho_h\,h_{i,t-1}+g_t+\xi_{it}, & \xi_{it}&\sim N(0,\sigma_\xi^2), \label{eq:h}\\
g_t &= \rho_g\,g_{t-1}+\nu_t, & \nu_t&\sim N(0,\sigma_g^2). \label{eq:g}
\end{align}
The global component $g_t$ is common to all countries and represents world-wide temperature fluctuations; it is the only source of cross-sectional dependence in the design. Endogeneity of emissions is induced through a shared innovation,
\begin{equation}\label{eq:endog}
\varepsilon_{it}=u_{it}+\lambda\eta_{it},\qquad u_{it}\sim N(0,\sigma_u^2),\qquad \lambda\in\{0,\,0.3,\,0.6,\,0.9\}.
\end{equation}
Hence $\operatorname{Corr}(\varepsilon_{it},\eta_{it})=\lambda\sigma_\eta/(\sigma_u^2+\lambda^2\sigma_\eta^2)^{1/2}$, which equals 0, 0.514, 0.768 and 0.874 for the four settings. Because $\eta_{it}$ is serially independent, $\varepsilon_{it}$ is serially uncorrelated, $q_{it}$ is contemporaneously endogenous, $q_{i,t-1}$ is predetermined, and precipitation and temperature are strictly exogenous. Table~\ref{tab:design} lists every parameter value.

\paragraph{Initial conditions.} Each series is started at its country-specific steady state $B=50$ periods before the sample, with $q_{i,-B}=m^q_i$, $p_{i,-B}=m^p_i$, $h_{i,-B}=g_{-B}=0$ and $y_{i,-B}=(\mu_i+\beta_1 m^q_i+\beta_2 m^p_i)/(1-\alpha)$, and the first $B$ periods are discarded. Because the start coincides with each country's long-run mean, the deviations of $y_{it}$ and $q_{it}$ from their country means have expectation zero for every $t$; the process is mean-stationary, and lagged differences are uncorrelated with $\mu_i$. This is the initial-condition restriction that System-GMM requires \citep{Blundell1998}, so the design represents the case most favourable to that estimator.

\paragraph{Random numbers.} Draws use R's default Mersenne-Twister generator with the R~$\ge$~3.6 sampling method. Replication $r=0,\dots,999$ of setting $\lambda_k$ ($k=0,\dots,3$) is generated after \texttt{set.seed(20261030 + 100000*k + r)}, so any replication can be regenerated on its own, independently of the number of processor cores used.

\begin{table}[htbp]
\centering
\begin{threeparttable}
\caption{Parameter values of the data-generating process}\label{tab:design}
\small
\begin{tabular}{@{}lll@{}}
\toprule
Component & Parameter & Value \\
\midrule
Panel & Initial countries $N_0$; years $T$; burn-in $B$; replications $R$ & 135; 50; 50; 1{,}000 per $\lambda$ \\
GDP equation \eqref{eq:model} & $\alpha,\ \beta_1,\ \beta_2,\ \beta_3$ & 0.70, 0.70, 0.003, $-0.015$ \\
 & $\sigma_u$ (sd of $u_{it}$) & 0.05 \\
Country effects \eqref{eq:mu}--\eqref{eq:mp} & $\sigma_\mu$ & 0.30 \\
 & $\bar m_q,\ \sigma_{mq},\ \varrho$ & 5.55, 0.30, 0.60 \\
 & $\bar m_p,\ \sigma_{mp}$ & $\ln 870$, 0.25 \\
\COtwo{} \eqref{eq:q} & $\rho_q,\ \sigma_\eta$ & 0.80, 0.10 \\
Precipitation \eqref{eq:p} & $\rho_p,\ \sigma_\zeta$ & 0.30, 0.25 \\
Temperature \eqref{eq:h}--\eqref{eq:g} & $\rho_h,\ \sigma_\xi;\ \rho_g,\ \sigma_g$ & 0.60, 0.70; 0.50, 0.30 \\
Endogeneity \eqref{eq:endog} & $\lambda$ & 0, 0.3, 0.6, 0.9 \\
Imbalance \eqref{eq:obs} & Max.\ entry offset; deletion prob.; min.\ observed years & 30; 0.08; 20 \\
\bottomrule
\end{tabular}
\begin{tablenotes}\footnotesize
\item Note: All innovations are independent Gaussian. The implied unconditional standard deviation of the temperature anomaly is approximately 1\,\textdegree C.
\end{tablenotes}
\end{threeparttable}
\end{table}

\subsection{Panel imbalance}

Imbalance is imposed after the complete panel has been generated. Each country receives an entry offset $e_i$ drawn uniformly from $\{0,\dots,30\}$, and each post-entry observation is retained independently with probability $0.92$:
\begin{equation}\label{eq:obs}
r_{it}=\mathbb{1}\{t>e_i\}\,b_{it},\qquad b_{it}\sim\text{Bernoulli}(0.92).
\end{equation}
Countries with fewer than 20 observed years are dropped, so 135 is the initial rather than the retained number of countries. Missingness is completely at random conditional on entry. Lags are defined in calendar time: $\Delta y_{it}$ exists only if years $t$ and $t-1$ are both observed, and a row enters an estimating equation only when all of its variables are observed.

\subsection{Estimators and instrument sets}\label{sec:est}

\paragraph{Fixed effects.} OLS on the within-transformed Eq.~\eqref{eq:model}, using rows with $y_{it}$ and $y_{i,t-1}$ observed and demeaning over each country's usable rows. Standard errors are clustered by country.

\paragraph{Instrumental variables.} The Anderson--Hsiao estimator applies two-stage least squares to the first-differenced equation
\begin{equation}\label{eq:fd}
\Delta y_{it}=\alpha\Delta y_{i,t-1}+\beta_1\Delta q_{it}+\beta_2\Delta p_{it}+\beta_3\Delta h_{it}+\Delta\varepsilon_{it},
\end{equation}
with no intercept. $\Delta y_{i,t-1}$ and $\Delta q_{it}$ are endogenous, the excluded instruments are $y_{i,t-2}$ and $q_{i,t-2}$, and $\Delta p_{it}$, $\Delta h_{it}$ instrument themselves. The model is exactly identified, so no overidentification test is available, and the estimator has no finite moments \citep{Stock2005}; its mean-based performance measures must therefore be read alongside median-based ones. Standard errors are clustered by country.

\paragraph{Difference-GMM.} Eq.~\eqref{eq:fd} is estimated with the moment conditions
\begin{equation}\label{eq:dgmm}
\E[y_{i,t-s}\Delta\varepsilon_{it}]=0,\quad \E[q_{i,t-s}\Delta\varepsilon_{it}]=0,\quad s\in\{2,3\};\qquad \E[\Delta p_{it}\Delta\varepsilon_{it}]=\E[\Delta h_{it}\Delta\varepsilon_{it}]=0.
\end{equation}
Lag 1 of $q$ is excluded because $q_{i,t-1}$ contains $\eta_{i,t-1}$, which is correlated with $\varepsilon_{i,t-1}$ and hence with $\Delta\varepsilon_{it}$. Restricting lags to 2--3 and collapsing the instrument matrix \citep{Roodman2009b} gives \textbf{6 instruments} for 4 parameters (Hansen df $=2$) in every replication; a lag whose calendar year is unobserved is set to zero, as is standard for collapsed instruments.

\paragraph{System-GMM.} The differenced equations use the instruments of Eq.~\eqref{eq:dgmm}. The level equations~\eqref{eq:model}, with an intercept, add
\begin{equation}\label{eq:sgmm}
\E[\Delta y_{i,t-1}(\mu_i+\varepsilon_{it})]=0,\quad \E[\Delta q_{i,t-1}(\mu_i+\varepsilon_{it})]=0,\quad \E[p_{it}(\mu_i+\varepsilon_{it})]=\E[h_{it}(\mu_i+\varepsilon_{it})]=0,
\end{equation}
plus the constant. The first two conditions hold because of mean stationarity (Section~\ref{sec:dgp}); the last two hold because $p_{it}$ and $h_{it}$ are independent of $\mu_i$ and $\varepsilon_{it}$ in this design. This gives \textbf{11 instruments} for 5 parameters (Hansen df $=6$) in every replication.

\paragraph{GMM estimation.} Both GMM estimators use two-step efficient weighting. The first step uses the $H$ matrix implied by first-differencing serially uncorrelated errors for the differenced rows and the identity for level rows. Two-step standard errors use the \citet{Windmeijer2005} correction.

\subsection{Specification diagnostics}

For each GMM fit we compute the two-step \citet{Hansen1982} $J$ test and, because two-step $J$ can be size-distorted when the first-step weight is inefficient, the $J$ statistic from iterated GMM \citep{Hansen1996}. We also compute the Arellano--Bond AR(1) and AR(2) tests on differenced residuals, including the parameter-estimation correction terms. For FE residuals we compute the \citet{Pesaran2015} CD test, using pairwise overlapping years ($\ge3$), and the \citet{Wooldridge2010} serial-correlation test \citep{Drukker2003}. For IV we report cluster-robust first-stage $F$ statistics for the two excluded instruments. Because the true disturbances are known in a simulation, we also verify each moment condition directly. In every replication we compute the sample correlation between each instrument and the true transformed disturbance: $\Delta\varepsilon_{it}$ for differenced instruments and $\mu_i+\varepsilon_{it}$ for level instruments.

\subsection{Performance measures}

Let $\hat\theta^{(r)}$ be the estimate of coefficient $\theta_{0}$ in replication $r$ and $\bar\theta$ its mean over the $R_s$ successful fits. Following \citet{Morris2019}, we report
\begin{equation}\label{eq:metrics}
\begin{aligned}
\text{Bias}&=\bar\theta-\theta_0,\qquad \text{RB}=100\,(\bar\theta-\theta_0)/\theta_0,\qquad
\text{SD}=\Big[\tfrac{1}{R_s-1}\textstyle\sum_r(\hat\theta^{(r)}-\bar\theta)^2\Big]^{1/2},\\
\text{RMSE}&=\Big[\tfrac{1}{R_s}\textstyle\sum_r(\hat\theta^{(r)}-\theta_0)^2\Big]^{1/2},\qquad
\text{MAE}=\operatorname{median}_r|\hat\theta^{(r)}-\theta_0|,
\end{aligned}
\end{equation}
together with the median bias and the coverage of nominal 95\% Wald intervals. The Monte Carlo standard error (MCSE) of the bias is $\text{SD}/\sqrt{R_s}$. A fit is counted as failed if a weighting or moment matrix is singular or an estimate or standard error is not finite; failed fits are reported and excluded from the measures.

\subsection{Software and reproducibility}

The simulation is implemented in R version 4.3.3 \citep{RCore2024}. All estimators and tests are coded in base R from their published formulae, so results do not depend on the defaults of any econometrics package; the \texttt{jsonlite} package \citep{Ooms2014} is used only to store output, and \texttt{parallel} to distribute replications across cores. The complete code, seeds and replication-level output are provided as Supplementary Material~S1. The full study of 4{,}000 replications runs in about two minutes on a two-core computer.

%% =====================================================================
\section{Results}\label{sec:results}

\subsection{Validation of the simulation}\label{sec:valid}

Table~\ref{tab:valid} summarises the samples and diagnostics across all replications. Every estimator produced a usable fit in all 4{,}000 replications, so the performance measures are based on $R_s=1{,}000$ throughout. On average 127.4 of the 135 countries were retained (range 117--135), providing about 4{,}200 observed country-years. Because differencing and lagging require consecutive observed years, the usable estimation rows fall to about 3{,}770 for FE and 3{,}375 for the differenced equation; System-GMM stacks about 7{,}145 differenced and level rows. The instrument count was fixed at 6 (Difference-GMM) and 11 (System-GMM) in every replication, well below the number of countries.

\begin{table}[htbp]
\centering
\begin{threeparttable}
\caption{Simulation validation: samples, instruments and replication-level diagnostics}\label{tab:valid}
\scriptsize
\setlength{\tabcolsep}{3pt}
\begin{tabular}{@{}p{5.3cm}cccc@{}}
\toprule
 & $\lambda=0$ & $\lambda=0.3$ & $\lambda=0.6$ & $\lambda=0.9$ \\
\midrule
%%TAB_VALID%%
\bottomrule
\end{tabular}
\begin{tablenotes}\footnotesize
\item Note: Rejection rates are shares of 1{,}000 replications rejecting at the 5\% level. The usable-row figures are means; Sys-GMM counts stacked differenced and level rows. The Wooldridge test assumes strictly exogenous regressors (Section~\ref{sec:recon}).
\end{tablenotes}
\end{threeparttable}
\end{table}

Table~\ref{tab:oracle} verifies the moment conditions directly. The mean correlation between every GMM instrument and the true disturbance is at most 0.0012 in absolute value, and the 95th percentile of its absolute value across replications never exceeds 0.03. This is the sampling noise expected of a zero correlation estimated from about 3{,}400 rows. The level instruments $\Delta y_{i,t-1}$ and $\Delta q_{i,t-1}$ are uncorrelated with $\mu_i+\varepsilon_{it}$, confirming that the initialisation delivers the mean-stationarity required by System-GMM. By contrast, the within-country correlation between $q_{it}$ and $\varepsilon_{it}$ rises from zero to 0.52 with $\lambda$, confirming that the endogeneity mechanism operates as intended. All instrument sets are therefore valid by construction and in sample, so the performance differences reported below arise from instrument strength and finite-sample behaviour, not from invalid moments.

\begin{table}[htbp]
\centering
\begin{threeparttable}
\caption{Direct check of the moment conditions: correlation of each instrument with the true disturbance}\label{tab:oracle}
\scriptsize
\setlength{\tabcolsep}{3.5pt}
\begin{tabular}{@{}lcccccccc@{}}
\toprule
 & \multicolumn{2}{c}{$\lambda=0$} & \multicolumn{2}{c}{$\lambda=0.3$} & \multicolumn{2}{c}{$\lambda=0.6$} & \multicolumn{2}{c}{$\lambda=0.9$} \\
\cmidrule(lr){2-3}\cmidrule(lr){4-5}\cmidrule(lr){6-7}\cmidrule(lr){8-9}
Instrument vs disturbance & Mean & P95$|r|$ & Mean & P95$|r|$ & Mean & P95$|r|$ & Mean & P95$|r|$ \\
\midrule
%%TAB_ORACLE%%
\bottomrule
\end{tabular}
\begin{tablenotes}\footnotesize
\item Note: Mean and 95th percentile of the absolute value of the sample correlation across 1{,}000 replications. Differenced-equation instruments are evaluated on non-zero (observed) entries.
\end{tablenotes}
\end{threeparttable}
\end{table}

The specification tests behave very differently from one another. Across replications, the AR(2) test rejected at close to its nominal 5\% rate for System-GMM in all settings and for Difference-GMM except at $\lambda=0.9$ (10.4\%). There, weak identification of the \COtwo{} coefficient distorts the residuals. The Difference-GMM Hansen test rejected at close to the nominal rate (3.5--5.2\%). The two-step Hansen test for System-GMM, however, rejected the valid instruments in 13.2--18.1\% of replications. The same moments evaluated by iterated GMM rejected in 4.5--7.8\%, close to the nominal level. The two-step over-rejection therefore reflects the inefficient first-step weight for the stacked system, not invalid instruments. Section~\ref{sec:recon} discusses the FE residual diagnostics.

\subsection{Illustrative single replication}\label{sec:single}

To show how the estimators appear to an applied researcher who observes a single panel, Tables~\ref{tab:desc}--\ref{tab:diag} report replication $r=0$ of the strong-endogeneity setting ($\lambda=0.9$), which retains 121 countries and 3{,}924 country-years. On the original scale (Table~\ref{tab:desc}), GDP is strongly right-skewed, emissions and precipitation are moderately right-skewed, and temperature anomalies are approximately symmetric about zero, which supports modelling the first three in logarithms. Pooled AR(1) coefficients (Table~\ref{tab:ar1}) are near unity for log GDP and log \COtwo{}. Because they mix within-country persistence with permanent cross-country differences, they exceed the true within-country autoregressive parameters (Table~\ref{tab:design}).

\begin{table}[htbp]
\centering
\begin{threeparttable}
\caption{Descriptive statistics of the illustrative replication ($\lambda=0.9$, $r=0$)}\label{tab:desc}
\small
\begin{tabular}{@{}lrrrr@{}}
\toprule
Statistic & GDP & \COtwo{} & Precipitation (mm) & Temp.\ anomaly (\textdegree C) \\
\midrule
%%TAB_DESC%%
\bottomrule
\end{tabular}
\begin{tablenotes}\footnotesize
\item Note: GDP, \COtwo{} and precipitation are reported before log transformation, in simulation units.
\end{tablenotes}
\end{threeparttable}
\end{table}

\begin{table}[htbp]
\centering
\begin{threeparttable}
\caption{Pooled AR(1) persistence coefficients in the illustrative replication}\label{tab:ar1}
\small
\begin{tabular}{@{}lcc@{}}
\toprule
Variable & AR(1) coefficient & Cluster-robust SE \\
\midrule
%%TAB_AR1%%
\bottomrule
\end{tabular}
\begin{tablenotes}\footnotesize
\item Note: Pooled OLS of each variable on its first lag and an intercept.
\end{tablenotes}
\end{threeparttable}
\end{table}

Table~\ref{tab:single} anticipates the Monte Carlo results. FE understates persistence (0.457) and doubles the \COtwo{} elasticity (1.445). IV recovers persistence (0.700), but its \COtwo{} estimate (0.424) has a standard error of 2.80, which leaves the elasticity effectively unidentified. Difference-GMM also recovers persistence (0.697) but overstates the \COtwo{} elasticity (1.110, standard error 0.894). System-GMM gives 0.772 and 0.542, with much smaller standard errors. The diagnostics in Table~\ref{tab:diag} would not alert a researcher to the FE bias. The IV first-stage $F$ statistics of 5.1 and 0.1 do reveal the weakness of its instruments, while the GMM Hansen and AR(2) tests do not reject.

\begin{table}[htbp]
\centering
\begin{threeparttable}
\caption{Coefficient estimates in the illustrative replication ($\lambda=0.9$, $r=0$)}\label{tab:single}
\small
\begin{tabular}{@{}lcccc@{}}
\toprule
 & FE & IV & Diff-GMM & Sys-GMM \\
\midrule
%%TAB_SINGLE%%
\bottomrule
\end{tabular}
\begin{tablenotes}\footnotesize
\item Note: Standard errors in parentheses: cluster-robust (FE, IV) and Windmeijer-corrected two-step (GMM). $^{***}p<0.01$, $^{**}p<0.05$, $^{*}p<0.10$.
\end{tablenotes}
\end{threeparttable}
\end{table}

\begin{table}[htbp]
\centering
\begin{threeparttable}
\caption{Specification diagnostics in the illustrative replication}\label{tab:diag}
\small
\begin{tabular}{@{}llrrr@{}}
\toprule
Model & Test & Statistic & df & $p$-value \\
\midrule
%%TAB_DIAG%%
\bottomrule
\end{tabular}
\begin{tablenotes}\footnotesize
\item Note: Diff-GMM uses 6 instruments and Sys-GMM 11. First-stage $F$ statistics are cluster-robust.
\end{tablenotes}
\end{threeparttable}
\end{table}

\subsection{Monte Carlo performance: persistence and \COtwo{} elasticity}\label{sec:mc}

Table~\ref{tab:perf} and Figure~\ref{fig:rmse} report the performance of the estimators for $\alpha$ and $\beta_1$. MCSEs of the bias are at most 0.004 for FE and System-GMM and at most 0.024 for Difference-GMM, so the differences emphasised below are well beyond simulation noise. The IV moments are dominated by outliers and are interpreted through the median measures.

\begin{table}[htbp]
\centering
\begin{threeparttable}
\caption{Monte Carlo performance for persistence ($\alpha=0.70$) and the \COtwo{} elasticity ($\beta_1=0.70$), $R=1{,}000$}\label{tab:perf}
\scriptsize
\setlength{\tabcolsep}{4pt}
\begin{tabular}{@{}llrrrrrrr@{}}
\toprule
Est. & Coef. & Bias & RB (\%) & SD & RMSE & Med.\ bias & MAE & Coverage \\
\midrule
%%TAB_PERF%%
\bottomrule
\end{tabular}
\begin{tablenotes}\footnotesize
\item Note: RB = relative bias; MAE = median absolute error; Coverage = share of nominal 95\% intervals containing the true value. Bold marks the lowest RMSE within each $\lambda$ and coefficient. IV has no finite moments (exact identification), so its bias, SD and RMSE are driven by outliers and are shown for completeness only.
\end{tablenotes}
\end{threeparttable}
\end{table}

\begin{figure}[htbp]
\centering
\begin{tikzpicture}
\begin{groupplot}[
  group style={group size=2 by 1, horizontal sep=1.8cm},
  width=0.46\textwidth, height=0.38\textwidth,
  xlabel={Endogeneity intensity $\lambda$},
  xtick={0,0.3,0.6,0.9}, xmin=-0.05, xmax=0.95,
  ymode=log, grid=major, grid style={gray!25},
  tick label style={font=\footnotesize}, label style={font=\footnotesize},
  title style={font=\small},
  legend style={font=\scriptsize, draw=none, fill=none},
  every axis plot/.append style={thick, mark size=2pt},
]
\nextgroupplot[title={(a) Persistence ($\alpha$)}, ylabel={RMSE (log scale)},
  legend to name=rmseleg, legend columns=4]
\addplot[color=black, mark=square*] coordinates %%FIG_A_FE%%;
\addlegendentry{FE}
\addplot[color=orange!85!black, mark=triangle*, dashed] coordinates %%FIG_A_IV%%;
\addlegendentry{IV}
\addplot[color=blue!70!black, mark=*] coordinates %%FIG_A_DG%%;
\addlegendentry{Diff-GMM}
\addplot[color=green!50!black, mark=diamond*] coordinates %%FIG_A_SG%%;
\addlegendentry{Sys-GMM}
\nextgroupplot[title={(b) \COtwo{} elasticity ($\beta_1$)}]
\addplot[color=black, mark=square*] coordinates %%FIG_B_FE%%;
\addplot[color=orange!85!black, mark=triangle*, dashed] coordinates %%FIG_B_IV%%;
\addplot[color=blue!70!black, mark=*] coordinates %%FIG_B_DG%%;
\addplot[color=green!50!black, mark=diamond*] coordinates %%FIG_B_SG%%;
\end{groupplot}
\end{tikzpicture}

\pgfplotslegendfromname{rmseleg}
\caption{RMSE of the persistence and \COtwo{} coefficients by estimator and endogeneity intensity (1{,}000 replications, log scale).}\label{fig:rmse}
\end{figure}

\paragraph{Exogenous baseline ($\lambda=0$).} With exogenous emissions FE is the most accurate estimator by a wide margin: its biases are below 1\% and its RMSEs (0.006 for $\alpha$, 0.008 for $\beta_1$) are less than one-fifth of those of System-GMM. The Nickell bias is small here because about 30 usable years are available per country, and because most of the within-country variation in $y_{i,t-1}$ comes from the strictly exogenous regressors rather than from past disturbances. The GMM estimators are consistent but pay a large efficiency cost, and System-GMM shows a small upward bias in $\alpha$ (2.7\%).

\paragraph{Endogenous emissions ($\lambda>0$).} The FE bias grows almost linearly with $\lambda$, from $-11.8\%$ to $-34.0\%$ for $\alpha$ and from $+31.5\%$ to $+105.8\%$ for $\beta_1$, and its intervals never cover the true values. Because the $\lambda=0$ bias is negligible, virtually all of this error is attributable to endogeneity rather than to the dynamic structure. IV is weakly identified: the median first-stage $F$ for $\Delta q_{it}$ falls from 5.2 to 1.1 (Table~\ref{tab:valid}), the median absolute error of $\hat\beta_1$ rises from 0.27 to 0.58, and occasional extreme estimates produce RMSEs of 7 to 12. Difference-GMM remains nearly unbiased for $\alpha$ when $\lambda\le0.6$ but is imprecise for $\beta_1$ (SD 0.38--0.75), and its bias for $\beta_1$ increases from 13\% to 74\%, with coverage falling to 81\% at $\lambda=0.9$. System-GMM has the lowest RMSE for both coefficients in every endogenous setting, with relative biases of $+2.2$ to $+2.8\%$ for $\alpha$ and $-1.6$ to $-2.8\%$ for $\beta_1$. Its RMSE is 59--63\% below that of Difference-GMM for $\alpha$ and 85--86\% below for $\beta_1$; at $\lambda=0.9$ its \COtwo{} RMSE (0.124) is 86.4\% below that of Difference-GMM (0.911) and 83.3\% below that of FE (0.741). Its Windmeijer-corrected intervals nevertheless under-cover, at 88--89\% for $\alpha$ and 92--93\% for $\beta_1$.

\subsection{Climatic coefficients}\label{sec:clim}

Table~\ref{tab:clim} reports the precipitation and temperature coefficients. All estimators except IV recover them with useful precision in this design, but the ranking differs from that for $\alpha$ and $\beta_1$. FE attains the lowest RMSE for $\beta_2$ in every setting and for $\beta_3$ when $\lambda\le0.3$, and ties with the GMM estimators at $\lambda=0.6$. Its error for $\beta_3$, however, becomes increasingly a matter of bias: it overstates the temperature effect by 22\% at $\lambda=0.9$, and its interval coverage falls to 26\%. Both GMM estimators are essentially unbiased for the climatic coefficients, with coverage at or slightly above the nominal level, Difference-GMM having the lowest $\beta_3$ RMSE at $\lambda=0.9$. IV is unbiased in median but far less precise. For strictly exogenous climatic regressors, therefore, FE is efficient but its inference becomes unreliable when another regressor is endogenous, whereas GMM trades some precision for valid inference.

\begin{table}[htbp]
\centering
\begin{threeparttable}
\caption{Monte Carlo performance for the climatic coefficients: bias, RMSE and coverage}\label{tab:clim}
\scriptsize
\setlength{\tabcolsep}{3pt}
\begin{tabular}{@{}llrrrrrrrrrrrr@{}}
\toprule
 & & \multicolumn{3}{c}{$\lambda=0$} & \multicolumn{3}{c}{$\lambda=0.3$} & \multicolumn{3}{c}{$\lambda=0.6$} & \multicolumn{3}{c}{$\lambda=0.9$} \\
\cmidrule(lr){3-5}\cmidrule(lr){6-8}\cmidrule(lr){9-11}\cmidrule(lr){12-14}
Coef. & Est. & Bias & RMSE & Cov. & Bias & RMSE & Cov. & Bias & RMSE & Cov. & Bias & RMSE & Cov. \\
\midrule
%%TAB_CLIM%%
\bottomrule
\end{tabular}
\begin{tablenotes}\footnotesize
\item Note: True values $\beta_2=0.003$ (precipitation elasticity) and $\beta_3=-0.015$ (temperature semi-elasticity). Bold marks the lowest RMSE within each $\lambda$ and coefficient. Cov.\ = coverage of nominal 95\% intervals.
\end{tablenotes}
\end{threeparttable}
\end{table}

%% =====================================================================
\section{Discussion}\label{sec:discussion}

\subsection{Sources of bias}

The $\lambda=0$ baseline separates two sources of bias that are usually conflated. In this design the dynamic-panel (Nickell) bias of FE is negligible, and the large FE biases observed for $\lambda>0$ come almost entirely from the endogeneity of emissions. Endogeneity contaminates both coefficients: FE overstates $\beta_1$ and, because $y_{i,t-1}$ and $q_{it}$ are strongly correlated, compensates by understating $\alpha$, so the implied long-run \COtwo{} elasticity is distorted on two fronts. The practical implication is that the case against FE in climate--economy models rests on the endogeneity of regressors such as emissions more than on the presence of a lagged dependent variable, at least when the time dimension is as long as here. Where emissions can be treated as exogenous, FE may be preferable to GMM.

The IV results show that a consistent estimator can be unusable in finite samples. With persistent $y$ and $q$, two-period-lagged levels carry little information about subsequent first differences, and the problem worsens as $\lambda$ rises because a larger share of the variation in $\Delta q_{it}$ comes from the endogenous innovation. Difference-GMM adds only one further lag and therefore inherits part of this weakness, which explains its imprecise and upward-biased \COtwo{} estimates. System-GMM overcomes it because lagged differences are informative about levels when series are persistent \citep{Blundell1998}.

\subsection{Reconciling the diagnostics}\label{sec:recon}

Specification tests from a single panel can mislead, and the replication-level results in Table~\ref{tab:valid} show why. The design contains no cross-sectional dependence in $\varepsilon_{it}$; the only common component is the global temperature factor $g_t$, which enters through a strictly exogenous regressor. When FE is correctly specified ($\lambda=0$), the CD test on its residuals rejects in 5.2\% of replications, at the nominal level. As $\lambda$ increases, the rejection rate rises to 31.5\%. This dependence is generated by misspecification, not by the design: the biased FE coefficients leave part of $y_{i,t-1}$ and $h_{it}$ in the residuals, and both contain the common factor $g_t$. Detecting cross-sectional dependence in FE residuals can therefore signal estimator bias rather than a feature of the data, and a Driscoll--Kraay-type correction of standard errors would not address it.

The Wooldridge test rejects in every replication, even at $\lambda=0$ where $\varepsilon_{it}$ is serially independent by construction. The test assumes strictly exogenous regressors \citep{Wooldridge2010,Drukker2003}. In a dynamic model the first-differenced OLS regression on which it is based is inconsistent, so its residual autocorrelation (median $-0.33$ to $-0.18$) departs from the $-0.5$ the test requires. The rejections therefore reflect applying the test outside its assumptions, not serial correlation in the data; the Arellano--Bond AR(2) test, which is valid in this setting, rejects at close to its nominal rate. Finally, the CD test does not reject in the illustrative replication ($p=0.241$), even though at $\lambda=0.9$ it rejects in almost a third of replications. A single replication cannot separate design, misspecification and chance, which is why the conclusions here rest on the replication-level evidence.

\subsection{Implications for applied research}

The results support System-GMM for estimating persistence and the effect of an endogenous, persistent regressor, but only under conditions resembling those simulated. Three of these conditions matter. First, the System-GMM level moments hold by construction here, and in empirical panels countries may be far from their steady states, in which case these moments fail and System-GMM can be more biased than Difference-GMM. Second, its two-step Hansen test over-rejected valid instruments, so researchers should report iterated-GMM or difference-in-Hansen statistics alongside it rather than discard a specification on a single two-step rejection. Third, its intervals under-covered, so inference on $\alpha$ and $\beta_1$ should be treated as somewhat optimistic. When emissions are plausibly exogenous, FE is more accurate. When interest centres on exogenous weather variables, FE is efficient but its intervals lose coverage if another regressor is endogenous, and GMM estimates of the climatic coefficients provide valid, if less precise, inference.

\subsection{Limitations}

The conclusions are conditional on the design. The study fixes $N$, $T$, the persistence of output and emissions, the variance ratio $\sigma_\mu/\sigma_u$ and a missing-completely-at-random imbalance mechanism, and it imposes mean-stationary initial conditions that favour System-GMM. It does not consider violations of that restriction, outcome-dependent missingness, cross-sectional dependence in $\varepsilon_{it}$, non-linear temperature effects \citep{Burke2015}, heterogeneous slopes, or bias-corrected FE estimators \citep{Bruno2005}. Each of these is a natural extension, and the supplied code is structured so that any design parameter can be varied.

%% =====================================================================
\section{Conclusion}\label{sec:conclusion}

This study compared four estimators of a dynamic climate--economy panel under persistence, endogenous emissions, common climatic shocks and panel imbalance, using a fully specified and reproducible Monte Carlo design with 1{,}000 replications at each of four endogeneity settings. With exogenous emissions, FE was the most accurate estimator and its dynamic-panel bias was negligible. Once emissions were endogenous, FE bias grew with endogeneity intensity, to $-34\%$ for persistence and $+106\%$ for the \COtwo{} elasticity at $\lambda=0.9$. IV was weakly identified, and Difference-GMM estimated the \COtwo{} elasticity imprecisely and with increasing upward bias. System-GMM attained the lowest RMSE for persistence and the \COtwo{} elasticity in every endogenous setting, but its intervals under-covered and its two-step Hansen test over-rejected valid instruments. For the climatic coefficients no estimator dominated: FE was most precise, while the GMM estimators gave unbiased estimates with valid intervals.

These results support System-GMM as the preferred estimator only when the regressor of interest is endogenous and persistent, and when the mean-stationarity restriction is defensible for the countries studied. In that case it should be used with a restricted, collapsed instrument set, with instrument counts, AR(2) tests and both two-step and iterated Hansen statistics reported. They do not establish that System-GMM is superior for all unbalanced climate panels. Where emissions can be treated as exogenous, or where initial conditions are far from steady state, FE or Difference-GMM may be preferable. Standard single-sample diagnostics did not reveal the FE bias, so estimator choice should rest on the identification argument rather than on passing tests.

%% =====================================================================
\section*{Declarations}

\paragraph{Data and code availability.} No empirical data were used; all data are simulated. The complete R code, the random-number seeds and the replication-level output needed to reproduce every table and figure are submitted with the manuscript as Supplementary Material~S1, and will be deposited in a public repository with a persistent identifier on publication.\verify{insert repository DOI if deposited before submission}

\paragraph{Funding.} This research received no specific grant from any funding agency in the public, commercial or not-for-profit sectors.\verify{confirm}

\paragraph{Competing interests.} The authors declare no competing interests.

\paragraph{Author contributions (CRediT).} \verify{assign roles to each author, e.g.\ E.\,Opoku: Conceptualisation, Supervision, Writing -- review \& editing; A.\,A.\,Adu: Methodology, Validation; H.\,Emmanuel: Formal analysis, Visualisation; O.\,Enock: Software, Data curation, Writing -- original draft. All authors read and approved the final manuscript.}

%% =====================================================================
\begin{thebibliography}{99}
\small
\setlength{\itemsep}{1pt}

\bibitem[Albarran and Arellano(2019)]{Albarran2019}
Albarran, P. and Arellano, M. (2019). GMM estimation from incomplete and rotating panels. \emph{Annals of Economics and Statistics}, 134, 5--42. \doi{10.15609/annaeconstat2009.134.0005}

\bibitem[Anderson and Hsiao(1981)]{Anderson1981}
Anderson, T.~W. and Hsiao, C. (1981). Estimation of dynamic models with error components. \emph{Journal of the American Statistical Association}, 76(375), 598--606. \doi{10.1080/01621459.1981.10477691}

\bibitem[Anderson and Hsiao(1982)]{Anderson1982}
Anderson, T.~W. and Hsiao, C. (1982). Formulation and estimation of dynamic models using panel data. \emph{Journal of Econometrics}, 18(1), 47--82. \doi{10.1016/0304-4076(82)90095-1}

\bibitem[Arellano and Bond(1991)]{Arellano1991}
Arellano, M. and Bond, S. (1991). Some tests of specification for panel data: Monte Carlo evidence and an application to employment equations. \emph{The Review of Economic Studies}, 58(2), 277--297. \doi{10.2307/2297968}

\bibitem[Arellano and Bover(1995)]{Arellano1995}
Arellano, M. and Bover, O. (1995). Another look at the instrumental variable estimation of error-components models. \emph{Journal of Econometrics}, 68(1), 29--51. \doi{10.1016/0304-4076(94)01642-D}

\bibitem[Atak et~al.(2011)]{Atak2011}
Atak, A., Linton, O. and Xiao, Z. (2011). A semiparametric panel model for unbalanced data with application to climate change in the United Kingdom. \emph{Journal of Econometrics}, 164(1), 92--115. \doi{10.1016/j.jeconom.2011.02.008}

\bibitem[Auffhammer et~al.(2013)]{Auffhammer2013}
Auffhammer, M., Hsiang, S.~M., Schlenker, W. and Sobel, A. (2013). Using weather data and climate model output in economic analyses of climate change. \emph{Review of Environmental Economics and Policy}, 7(2), 181--198. \doi{10.1093/reep/ret016}

\bibitem[Blundell and Bond(1998)]{Blundell1998}
Blundell, R. and Bond, S. (1998). Initial conditions and moment restrictions in dynamic panel data models. \emph{Journal of Econometrics}, 87(1), 115--143. \doi{10.1016/S0304-4076(98)00009-8}

\bibitem[Bruno(2005)]{Bruno2005}
Bruno, G.~S.~F. (2005). Approximating the bias of the LSDV estimator for dynamic unbalanced panel data models. \emph{Economics Letters}, 87(3), 361--366. \doi{10.1016/j.econlet.2005.01.005}

\bibitem[Bun and Windmeijer(2010)]{Bun2010}
Bun, M.~J.~G. and Windmeijer, F. (2010). The weak instrument problem of the system GMM estimator in dynamic panel data models. \emph{The Econometrics Journal}, 13(1), 95--126. \doi{10.1111/j.1368-423X.2009.00299.x}

\bibitem[Burke et~al.(2015)]{Burke2015}
Burke, M., Hsiang, S.~M. and Miguel, E. (2015). Global non-linear effect of temperature on economic production. \emph{Nature}, 527(7577), 235--239. \doi{10.1038/nature15725}

\bibitem[Dell et~al.(2012)]{Dell2012}
Dell, M., Jones, B.~F. and Olken, B.~A. (2012). Temperature shocks and economic growth: Evidence from the last half century. \emph{American Economic Journal: Macroeconomics}, 4(3), 66--95. \doi{10.1257/mac.4.3.66}

\bibitem[Drukker(2003)]{Drukker2003}
Drukker, D.~M. (2003). Testing for serial correlation in linear panel-data models. \emph{The Stata Journal}, 3(2), 168--177. \doi{10.1177/1536867X0300300206}

\bibitem[Gao et~al.(2024)]{Gao2024}
Gao, M., Fang, X., Ge, R., Fan, Y.-P. and Wang, Y. (2024). Multiple serial correlations in global air temperature anomaly time series. \emph{PLOS ONE}, 19(7), e0306694. \doi{10.1371/journal.pone.0306694}

\bibitem[Hansen(1982)]{Hansen1982}
Hansen, L.~P. (1982). Large sample properties of generalized method of moments estimators. \emph{Econometrica}, 50(4), 1029--1054. \doi{10.2307/1912775}

\bibitem[Hansen et~al.(1996)]{Hansen1996}
Hansen, L.~P., Heaton, J. and Yaron, A. (1996). Finite-sample properties of some alternative GMM estimators. \emph{Journal of Business \& Economic Statistics}, 14(3), 262--280. \doi{10.1080/07350015.1996.10524656}

\bibitem[Hayakawa(2007)]{Hayakawa2007}
Hayakawa, K. (2007). Small sample bias properties of the system GMM estimator in dynamic panel data models. \emph{Economics Letters}, 95(1), 32--38. \doi{10.1016/j.econlet.2006.09.011}

\bibitem[Hsiang(2016)]{Hsiang2016}
Hsiang, S. (2016). Climate econometrics. \emph{Annual Review of Resource Economics}, 8, 43--75. \doi{10.1146/annurev-resource-100815-095343}

\bibitem[Judson and Owen(1999)]{Judson1999}
Judson, R.~A. and Owen, A.~L. (1999). Estimating dynamic panel data models: A guide for macroeconomists. \emph{Economics Letters}, 65(1), 9--15. \doi{10.1016/S0165-1765(99)00130-5}

\bibitem[Kiviet(1995)]{Kiviet1995}
Kiviet, J.~F. (1995). On bias, inconsistency, and efficiency of various estimators in dynamic panel data models. \emph{Journal of Econometrics}, 68(1), 53--78. \doi{10.1016/0304-4076(94)01643-E}

\bibitem[Morris et~al.(2019)]{Morris2019}
Morris, T.~P., White, I.~R. and Crowther, M.~J. (2019). Using simulation studies to evaluate statistical methods. \emph{Statistics in Medicine}, 38(11), 2074--2102. \doi{10.1002/sim.8086}

\bibitem[Nickell(1981)]{Nickell1981}
Nickell, S. (1981). Biases in dynamic models with fixed effects. \emph{Econometrica}, 49(6), 1417--1426. \doi{10.2307/1911408}

\bibitem[Ooms(2014)]{Ooms2014}
Ooms, J. (2014). The jsonlite package: A practical and consistent mapping between JSON data and R objects. \emph{arXiv preprint}, arXiv:1403.2805. \url{https://arxiv.org/abs/1403.2805}

\bibitem[Pesaran(2015)]{Pesaran2015}
Pesaran, M.~H. (2015). Testing weak cross-sectional dependence in large panels. \emph{Econometric Reviews}, 34(6--10), 1089--1117. \doi{10.1080/07474938.2014.956623}

\bibitem[R Core Team(2024)]{RCore2024}
R Core Team (2024). \emph{R: A Language and Environment for Statistical Computing} (Version 4.3.3). R Foundation for Statistical Computing, Vienna, Austria. \url{https://www.R-project.org/}

\bibitem[Roodman(2009a)]{Roodman2009a}
Roodman, D. (2009a). How to do xtabond2: An introduction to difference and system GMM in Stata. \emph{The Stata Journal}, 9(1), 86--136. \doi{10.1177/1536867X0900900106}

\bibitem[Roodman(2009b)]{Roodman2009b}
Roodman, D. (2009b). A note on the theme of too many instruments. \emph{Oxford Bulletin of Economics and Statistics}, 71(1), 135--158. \doi{10.1111/j.1468-0084.2008.00542.x}

\bibitem[Stock and Yogo(2005)]{Stock2005}
Stock, J.~H. and Yogo, M. (2005). Testing for weak instruments in linear IV regression. In D.~W.~K. Andrews and J.~H. Stock (Eds.), \emph{Identification and Inference for Econometric Models: Essays in Honor of Thomas Rothenberg} (pp.\ 80--108). Cambridge University Press, Cambridge.

\bibitem[Ullah et~al.(2018)]{Ullah2018}
Ullah, S., Akhtar, P. and Zaefarian, G. (2018). Dealing with endogeneity bias: The generalized method of moments (GMM) for panel data. \emph{Industrial Marketing Management}, 71, 69--78. \doi{10.1016/j.indmarman.2017.11.010}

\bibitem[Windmeijer(2005)]{Windmeijer2005}
Windmeijer, F. (2005). A finite sample correction for the variance of linear efficient two-step GMM estimators. \emph{Journal of Econometrics}, 126(1), 25--51. \doi{10.1016/j.jeconom.2004.02.005}

\bibitem[Wooldridge(2010)]{Wooldridge2010}
Wooldridge, J.~M. (2010). \emph{Econometric Analysis of Cross Section and Panel Data} (2nd ed.). MIT Press, Cambridge, MA.

\end{thebibliography}

\end{document}
